\section{Discussion}
\label{sec:discussion}

\subsection{Key Findings}

\textbf{Complete Orthogonality Is Stronger Than Expected.} We found Jaccard = 0.0, not just $< 0.3$. Static tools analyze code text; tests execute behavior---fundamentally different properties with no overlap.

\textbf{Static Analysis Provides Dominant Coverage.} 32.5\% have static-only errors vs 0\% exec-only. Static catches issues earlier in the error cascade.

\subsection{Limitations}

\textbf{L1: Canonical Solutions Invalidate H-M2.} We analyzed EvalPlus canonical solutions, not LLM-generated code. This is the critical limitation of this work: canonical solutions are designed to pass tests, yielding only 0.3\% behavioral failure rate vs. the 40\% threshold. \textbf{The H-M2 hypothesis (execution detects behavioral errors in static-clean code) cannot be evaluated with this experimental design.} A proper test requires generating code via LLM API calls where semantic errors naturally occur. The behavioral rate finding should not be interpreted as evidence about real LLM code behavior.

\textbf{L2: Combined Improvement Untested.} H-M3/H-M4 (combined $>$ max) not executed. The theoretical basis is established; empirical verification is future work.

\textbf{L3: Single Benchmark Family.} Results are for Python on HumanEval+/MBPP+. Cross-language generalization is future work.

\subsection{Broader Impact}

This research supports improved code generation systems by providing evidence for feedback composition. No negative societal impacts identified.
